\section{Related Work}
\label{sec:related_work}

Our proposed \textbf{HyperMerge} resides at the intersection of three highly active domains in machine learning: Parameter-Efficient Fine-Tuning (PEFT), Model Merging/Ensembling, and Hyperbolic Deep Learning.

\subsection{Parameter-Efficient Fine-Tuning (PEFT)}
To adapt massive foundation models without the prohibitive memory and storage requirements of full-parameter fine-tuning, PEFT has become the dominant adaptation strategy. Popular approaches include adapter layers \cite{houlsby2019parameter}, prefix-tuning \cite{li2021prefix}, prompt-tuning \cite{lester2021power}, and Low-Rank Adaptation (LoRA) \cite{hu2021lora}. Among these, LoRA has gained widespread adoption due to its lack of additional inference latency (as adapter weights can be folded back into the base model weights upon deployment) and its highly effective low-rank factorization of weight updates. Dynamic serving frameworks such as S-LoRA \cite{sheng2023slora} and Punica \cite{chen2023punica} serve thousands of concurrent LoRA adapters, but they require heavy system-level scheduling, memory-pinning, and paging mechanisms, which are extremely difficult to deploy on edge/CPU-bound platforms.

\subsection{Static Model Merging and Ensembling}
Model merging aims to fuse multiple independent task-specific models into a single parameter set without retraining. Weight averaging, or model soups \cite{wortsman2022model}, averages weights of models fine-tuned on the same task to improve generalization. When merging models fine-tuned on different tasks, task arithmetic \cite{ilharco2022editing} adds task vectors (differences between fine-tuned and pre-trained weights) to the pre-trained base model. To resolve the parameter conflicts and sign mismatches that degrade task arithmetic, TIES-merging \cite{yadav2023ties} trims small magnitude parameters and resolves sign conflicts. DARE \cite{yu2023dare} prunes redundant parameters using a delta-weight drop-and-rescale strategy. Additionally, Git Re-Basin \cite{ainsworth2022git} and ZipIt! \cite{stoica2023zipit} permute neurons to align weight spaces before merging. While these methods are powerful, their static nature makes them fundamentally unsuited for processing heterogeneous multi-task streams in real-time, resulting in performance collapse (to Uniform merging) when processing conflicting inputs.

\subsection{Dynamic Test-Time Adaptation and Routing}
To overcome the limitations of static merging, dynamic, test-time ensembling methods have emerged. AdaMerging \cite{pfsr2025} dynamically optimizes merging coefficients at test-time via online gradient descent on incoming batches, but it introduces a severe $26\times$ latency penalty due to backward passes. Parameter-Free Subspace Routing (PFSR) \cite{pfsr2025} computes closed-form task projections in a single forward pass, but under heterogeneous batches, it collapses to Uniform merging unless paired with Micro-Batch Homogenization (MBH) \cite{mbh2025}. However, MBH requires grouping similar samples, which introduces $O(K)$ sequential backbone passes and severely degrades latency on edge devices. Sample-wise Activation Blending (SABLE) performs dynamic activation-space blending to eliminate MBH's sequential bottlenecks but suffers from sub-optimal ensembling quality under overlapping task boundaries. The state-of-the-art Euclidean method, Single-Pass Centroid Alignment (SPS-ZCA), uses centroid alignment to route inputs in a single forward pass but is fundamentally constrained by flat coordinate systems that suffer from representation crowding and inter-task cross-talk near the origin.

\subsection{Hyperbolic Deep Learning}
Hyperbolic geometry has emerged as a powerful tool for learning representations of data with hierarchical or power-law structures, such as social networks, word taxonomies, and biological trees. Nickel and Kiela \cite{nickel2017poincare} demonstrated that Poincaré embeddings can represent complex hierarchical structures with extremely low dimensions and low distortion. Ganea et al.\ \cite{ganea2018hyperbolic} introduced Hyperbolic Neural Networks (HNNs), formulating fundamental neural network layers (e.g., multinomial logistic regression, pointwise feed-forward) using Möbius gyrovector spaces in the Poincaré Ball. Chami et al.\ \cite{chami2019hyperbolic} extended this to graph neural networks, showing that negative curvature is highly effective for processing hierarchical graph structures. While hyperbolic representations have been applied to graph learning, computer vision, and NLP, \textbf{HyperMerge} is the first work to introduce hyperbolic geometry and non-linear Möbius algebra to the domain of modular deep learning and dynamic model ensembling, breaking the traditional flat Euclidean barrier.
